\clearpage
\VerdierPhysicalPageAnchor{0162}%
\begingroup\small\setlength{\parskip}{0pt}\setlength{\emergencystretch}{1.5em}
\VerdierSemanticAnchor{ch4}
\section{Objets spectraux.}

Nous reprenons, pour l'adapter aux catégories triangulées, l'exposition des
suites spectrales de [1], chap. XV, \S 7.

\VerdierSemanticAnchor{ch4-4-1}
\subsection{Définition des objets spectraux.}

\VerdierSemanticAnchor{ch4-4-1-1}
\medskip\noindent\textbf{4.1.1.}\quad
Soit $\mathcal J$ une catégorie et, pour tout $n\in\mathbb N$, soit
$\mathcal{F}\ell_n(\mathcal J)$ la catégorie des suites de $n$ morphismes
composables de $\mathcal J$ (chap. I, 3.3.6). Nous utiliserons trois
foncteurs notés
$s_i\colon\mathcal{F}\ell_2(\mathcal J)\to\mathcal{F}\ell_1(\mathcal J)$,
$1\leq i\leq3$ :
\VerdierSemanticAnchor{ch4-4-1-1-1}
\[
\left\{
\begin{aligned}
&\bigl(\mathord\bullet\xrightarrow f\mathord\bullet\xrightarrow g\mathord\bullet\bigr)
   \in\operatorname{Ob}\mathcal{F}\ell_2(\mathcal J),\\
&s_1\bigl(\mathord\bullet\xrightarrow f\mathord\bullet\xrightarrow g\mathord\bullet\bigr)
   =\bigl(\mathord\bullet\xrightarrow f\mathord\bullet\bigr),\\
&s_2\bigl(\mathord\bullet\xrightarrow f\mathord\bullet\xrightarrow g\mathord\bullet\bigr)
   =\bigl(\mathord\bullet\xrightarrow{gf}\mathord\bullet\bigr),\\
&s_3\bigl(\mathord\bullet\xrightarrow f\mathord\bullet\xrightarrow g\mathord\bullet\bigr)
   =\bigl(\mathord\bullet\xrightarrow g\mathord\bullet\bigr).
\end{aligned}
\right.
\tag{4.1.1.1}
\]
et deux morphismes de foncteurs :
\VerdierSemanticAnchor{ch4-4-1-1-2}
\[
\left\{
\begin{aligned}
&u_1\colon s_1\longrightarrow s_2,
&&u_1(f,g)=
\begin{tikzcd}[ampersand replacement=\&,column sep=2.8em,row sep=1.9em,baseline=(current bounding box.center)]
{}\arrow[r,"f"]\arrow[d,"\operatorname{id}"']\&{}\arrow[d,"g"]\\
{}\arrow[r,"gf"']\&{}
\end{tikzcd},\\[2pt]
&u_2\colon s_2\longrightarrow s_3,
&&u_2(f,g)=
\begin{tikzcd}[ampersand replacement=\&,column sep=2.8em,row sep=1.9em,baseline=(current bounding box.center)]
{}\arrow[r,"gf"]\arrow[d,"f"']\&{}\arrow[d,"\operatorname{id}"]\\
{}\arrow[r,"g"']\&{}
\end{tikzcd}.
\end{aligned}
\right.
\tag{4.1.1.2}
\]

\VerdierSemanticAnchor{ch4-4-1-2}
\medskip\noindent\textbf{4.1.2.}\quad
Soient $\mathcal D$ une catégorie triangulée et $\mathcal J$ une catégorie.
Un \emph{objet spectral de type $\mathcal J$ à valeurs dans $\mathcal D$}
est un objet constitué par deux foncteurs :

a) Un foncteur
$X\colon\mathcal{F}\ell_1(\mathcal J)\longrightarrow\mathcal D$ :
\[
\bigl(\mathord\bullet\xrightarrow f\mathord\bullet\bigr)\longmapsto X_f.
\]
\endgroup
